\section{Conclusion}
\label{sec:conclusion}

We presented a direction-optimizing BFS that is wait-free under TSO, given that a thread can allocate memory in a bounded number of its own steps: every vertex receives its BFS distance, final once any call returns, and the work of a top-down level depends only on that level and on $T$. A variant also bounds its memory by the largest levels, plus a constant per level. On the five standard graphs, our implementation examined at most 30\% more neighbors than a sequential search at 112 threads, its best observed times were below those of GAPBS, GBBS and PASGAL, and with a thread stalled at every level it lost little, while GAPBS, which waits at every level, lost about one stall per level.

Since every thread takes part in every level, the method is weakest on graphs with many levels of few vertices, where blocking implementations that leave such a level to one thread are faster; and the reset before a search, outside the method, still waits for every thread. The main open questions are running ahead of small levels without changing a published distance, cheap orders with provable contention bounds, and reclaiming and resetting the state of a search without waiting for a stopped thread; others concern direction choices at high thread counts, hubs, and weaker memory models. Finally, nothing in the structure of the method, a list of levels whose per-thread buckets any thread can finish, is specific to BFS, and it could make other frontier-based graph algorithms wait-free, such as connectivity, or single-source shortest paths, whose buckets of tentative distances play the role of levels.
